% 从旧数学附录迁出的带修莫队说明；不新增算法。
\section{带修莫队：时间轴与区间状态}\label{knowledge-mo-time}
\index{带修莫队}
先离线读完操作，记录每个单点修改的下标、修改前值和修改后值。
询问的时间 $t$ 表示前 $t$ 次修改已执行；输入中的询问本身不推进时间。
同一位置的第二次修改，其旧值必须来自第一次修改后的数组，不能始终使用初始值。
需要离散化时，应将全部初值与修改后值一同纳入。

调度器保持三个指针 $(l,r,t)$，回调统计恰好对应当前版本的半开区间 $[l,r)$。
前进修改 $i$ 时写入 after，回退时写入 before。
若修改点在当前区间内，必须先删除旧贡献、改值、再加入新贡献；区间外只改数组。
随后移动区间端点，最后按原询问编号写答案。
扩展区间先于收缩，支持空区间和互不相交区间。
整个 run 结束后状态停留在排序后的最后一条询问；再次调用前须重置原数组和统计量。

左端点分块、右端点分块，再按时间排序，相邻块采用蛇形顺序。
时间关键字是真实修改次数，不能把询问编号当作时间维度。
令块长为 $B$，询问数为 $Q$，修改数为 $U$，则常数时间回调下的移动量上界为
\[
O\left(QB+\frac{n^2}{B}+U\min\left(Q,(\lfloor n/B\rfloor+1)^2\right)\right).
\]
另有 $O(Q\log Q)$ 排序与 $O(Q)$ 调度空间。
三者同阶时取 $B\approx n^{2/3}$，得到 $O(n^{5/3})$；规模悬殊时可手动调整块长。
回调需要区分左右端增删时应拆分四个方向，当前简版适合不同数、频次矩等统计。

\href{https://www.luogu.com.cn/problem/P1903}{P1903 数颜色} 驱动以颜色频次数组维护不同数，
真实操作顺序建立修改前后值，随后交给调度器。
本地快照测试逐项检查每个时间、每个区间及修改回退状态，另以顺序 set 查询核对完整驱动。
最大规模验证 133333 个元素和操作，并覆盖仅修改、仅询问与重复修改同一位置。
最终简短命名版本在 \href{https://www.luogu.com.cn/record/297653743}{P1903 记录 297653743}
通过全部 13 个测试，最慢单点 546 ms、峰值内存 14.65 MB。
此记录使用默认块长与不同颜色数统计；自定义块长、空区间及任意登记时间由本地测试覆盖。
最快提交榜的排名及总数尚未核验。

\newpage
